\begin{theorem}[G-18: Every large stretch in every strict four-row direction]\label{G-18}
Let $n>m\geq3$, $a=n-m$, and let $\nu\vdash ns$, $s\geq1$,
have at most four rows. Assume
\[
 \nu_4>0,\qquad \nu_1>as.
\]
Use $\eta,r_\eta,b,L,p,\beta,T,E$ from G-17, defined by the same
formulas even when its dimension inequality is not assumed. Put
\[
 c_m=\dim\mathcal P_{m,4},\qquad
 u=2n^2+r_\eta-13,\qquad v=c_m+r_\eta-12.
\]
The following conclusions hold in original coordinate degree $sj$.

First, both representations occur in every sufficiently large
integer stretch, with the uniform displayed bound
\begin{equation}\label{G18:occurrence}
 m_{j\nu}(\ell^a\operatorname{per}_m)\geq1,
 \qquad m_{j\nu}(\det_n)\geq1
 \qquad(j\geq\beta^2).
\end{equation}
This conclusion needs no dimension comparison between the two cones.

If $v>u$, equivalently $c_m+3>2n^2+2$, set
$Q_D=(pE)^{u-1}$. Then
\begin{equation}\label{G18:everylargepositive}
 m_{j\nu}(\ell^a\operatorname{per}_m)>m_{j\nu}(\det_n)
 \qquad\text{for every integer }j\geq E(4uQ_D+1).
\end{equation}
Thus every affirmative direction of G-17, including the rectangular
ones, gives a gap at every sufficiently large integer stretch.

There is also an effective converse range. Set
\[
 \bar v=4m^2-2m+r_\eta-10,\qquad Q_P=(pE)^{\bar v-1}.
\]
If
\begin{equation}\label{G18:reversecondition}
 2n^2>4m^2-2m+3,
\end{equation}
then
\begin{equation}\label{G18:everylargereverse}
 m_{j\nu}(\det_n)>m_{j\nu}(\ell^a\operatorname{per}_m)
 \qquad\text{for every integer }j\geq E(4\bar vQ_P+1).
\end{equation}
Condition \eqref{G18:reversecondition} holds, in particular,
whenever $n\geq\sqrt2\,m$.

More precisely, for every such fixed direction,
there is a constant $\kappa_\nu>0$ such that
\begin{equation}\label{G18:ratio}
 \frac{m_{j\nu}(\ell^a\operatorname{per}_m)}
      {m_{j\nu}(\det_n)}
  =\kappa_\nu j^{\,c_m-2n^2+1}(1+O(j^{-1}))
  \qquad(j\longrightarrow\infty)
\end{equation}
through all positive integers. When the exponent is zero, the
dimension comparison alone does not determine the sign of the gap.

For every fixed partition direction $\nu\vdash ns$ with at most
four rows, including those outside the two strict inequalities,
there is no occurrence obstruction at any sufficiently large
stretch. Here an occurrence obstruction means positive
padded-permanent multiplicity and zero determinant multiplicity.
\end{theorem}

\paragraph{Sources and the extra step.}
The enhanced coordinate rings and the initial stable tensor are
proved in G-17. We additionally use separation of distinct closed
orbits by invariants: Mukai (2003), Theorem~5.3, printed pp.~161--162,
and Corollary~5.4, p.~162. Openness of the stable locus is
Proposition~5.15, p.~166. The new step removes the finite projective
stabilizer by moving the padding line, then makes the resulting
degree-semigroup argument effective. The original coordinate degree
is $sj$; in a rectangular direction $ns=4\nu_4$ already incorporates
the necessary central congruence.

\begin{proof}
We retain the notation $X_{Y_D},X_{Y_P},W,A_D,A_P$ of G-17.
Their construction, common stable tensor, generator bound, and
dimensions $u,v$ do not require the inequality $v>u$.

\paragraph{1. A common stable tensor with no nontrivial scalar orbit symmetry.}
Write the balanced padded Fermat form of G-17 as $x_0^a f$, where
$f=\sum_{i=0}^3\gamma_i x_i^m$ and every $\gamma_i$ is nonzero.
The degree-$m$ hypersurface $f=0$ is smooth: its partial derivatives
have no common projective zero. It is irreducible, since two
positive-degree projective hypersurface components in $\mathbb P^3$
intersect, and their intersection would be singular in the product.
Its projective stabilizer $H\subseteq\mathrm{PGL}_4$ is finite.
Indeed, the equation $A\cdot f=cf$ forces all off-diagonal entries
of $A$ to vanish by the distinct monomials $x_i^{m-1}x_j$,
$i\ne j$, because $m\geq3$; its diagonal entries are then equal.
Thus the projective stabilizer has zero Lie algebra, so in
characteristic zero it is finite.

Keep $f$ and the auxiliary vector $v_\eta$ fixed, and vary a nonzero
padding linear form $\ell$. For every such $\ell$, the tensor
\[
 w_\ell=(\ell^a f)^{\otimes s}\otimes v_\eta
\]
belongs to both enhanced cones. This follows from the same binary
factorization and cyclic permanent/determinant representations of
$f$ used in G-17, followed by the independent padding block.
The inverse image of the stable locus of either cone under this
regular map is a nonempty open set, because it contains $\ell=x_0$.
Stability is unchanged by a nonzero scalar, so these sets descend
to nonempty open subsets of the projective space of linear forms.

For every nonidentity $h\in H$, its fixed locus in this projective
space is proper and closed: a linear transformation fixing every
line of the dual representation is scalar. The intersection of
the two stable opens with the complements of these finitely many
fixed loci is nonempty, by irreducibility of projective space.
Choose $[\ell]$ in this intersection. Because $a>0$ and $m\geq3$,
unique factorization implies that every projective automorphism of
$F=\ell^a f$ preserves the unique linear factor $\ell$ and the
unique irreducible factor $f$ of degree greater than one separately.
Its image therefore lies in $H$ and fixes $[\ell]$, so it is the
identity in $\mathrm{PGL}_4$.

Set $w=w_\ell$. A projective stabilizer of this decomposable tensor
preserves the line of $F^{\otimes s}$ and hence the line of $F$:
the Veronese map is injective on projective points. Thus its
projective stabilizer in $\mathrm{SL}_4$ is exactly the center
$\mu_4$. This center acts
on $W$ by the exponent
\[
 ns-|\eta|=4\nu_4,
\]
and hence acts trivially. It follows that
\begin{equation}\label{G18:noscaling}
 gw=cw,\quad g\in\mathrm{SL}_4,\ c\in\mathbb C^*
 \quad\Longrightarrow\quad c=1.
\end{equation}
The orbits of $w$ and of every $cw$ are closed: stability holds,
and scalar multiplication commutes with the group action.

\paragraph{2. Common nonzero invariant degrees have greatest common divisor one.}
Choose homogeneous generators $g_1,\ldots,g_t$ of
$\mathbb C[W]^{\mathrm{SL}_4}$, with positive degrees
$e_i\leq\beta$, as in G-17. Let
\[
 J=\{i:g_i(w)\ne0\}.
\]
This set is nonempty. Otherwise all positive-degree invariants
would vanish at $w$, so the distinct closed orbits of $w$ and zero
would fail to be separated.

Suppose the integer $q=\gcd\{e_i:i\in J\}$ exceeded one, and
choose a primitive $q$th root $c\ne1$. Homogeneity gives
$g_i(cw)=g_i(w)$ for every generator: for $i\in J$ use
$c^{e_i}=1$, and for $i\notin J$ both values are zero. The same
equality therefore holds for every invariant polynomial. Closed
orbit separation makes the two closed orbits equal, contrary to
\eqref{G18:noscaling}. Thus
\begin{equation}\label{G18:gcd}
 \gcd\{e_i:i\in J\}=1.
\end{equation}
Every product of these generators is nonzero at $w$, and consequently
has nonzero image in both $A_D$ and $A_P$.

Here is the required elementary numerical-semigroup bound. If
positive integers $e_i\leq\beta$ have greatest common divisor one,
their nonnegative semigroup contains every integer $N\geq\beta^2$.
Let $a_0$ be the least generator. In the directed graph on
$\mathbb Z/a_0\mathbb Z$ with edges adding the $e_i$, every residue
is reachable: a subsemigroup of a finite group is a subgroup, and
the greatest-common-divisor condition makes this subgroup the whole
group. A shortest path to each residue has no repeated vertex and
uses at most $a_0-1$ edges. Its integer weight $c_r$ therefore
satisfies $c_r\leq(a_0-1)\beta<\beta^2$. For $N\geq\beta^2$,
choose the representative of its residue; $N-c_r$ is a nonnegative
multiple of $a_0$. This proves the assertion, including $a_0=1$.

Applying this to \eqref{G18:gcd}, for every $N\geq\beta^2$ there
is a homogeneous ambient invariant $h_N$ of degree $N$ nonzero
at $w$. Its restrictions to both cones are nonzero. The exact
single-label identity of G-17 now proves \eqref{G18:occurrence}.

\paragraph{3. Transfer from the Veronese progression to every residue.}
We have $T\geq\beta$, because the positive-degree polynomial
space through degree $\beta$ contains the powers of one variable.
Since $\beta\geq2$ and $\beta!\geq\beta$, it follows that
$E=4T\beta!\geq\beta^2$.

For $j=Ek+r$, $0\leq r<E$, multiplication by the common nonzero
invariants $h_{E+r}$ and $h_{2E-r}$ is injective in each of the
domains $A_D,A_P$. In particular, for either $A$ and every $k\geq1$,
\begin{equation}\label{G18:sandwich}
 \dim A_{E(k-1)}\leq\dim A_{Ek+r}\leq\dim A_{E(k+2)}.
\end{equation}
Only existence of these invariants is used; no algorithm for
computing a minimal generating set is being asserted.

Suppose first $v>u$. G-17's standard-Veronese dimension and
image-degree bounds give
\[
 \dim(A_P)_{E(k-1)}\geq\binom{k+u-1}{u},\qquad
 \dim(A_D)_{E(k+2)}\leq Q_D\binom{k+u+1}{u-1}.
\]
The exact ratio of the binomials is
\[
 \frac{\binom{k+u-1}{u}}{\binom{k+u+1}{u-1}}
 =\frac{k(k+1)(k+2)}{u(k+u)(k+u+1)}
 \geq\frac{k}{4u}\qquad(k\geq u+1).
\]
For $k\geq4uQ_D+1$ this is strictly greater than $Q_D$.
Combining with \eqref{G18:sandwich} proves
\eqref{G18:everylargepositive} in every residue class.

For the converse bound, R-13's row--column torus gives
$c_m\leq4m^2-2m+2$, and hence $v\leq\bar v$.
Condition \eqref{G18:reversecondition} says exactly $u>\bar v$.
The enhanced permanent cone has a dense homogeneous
parametrization of degree $p=ns+|\eta|$:
\[
 (Z,L,B)\longmapsto
 \bigl(L^a\operatorname{per}_m(Z_0x_0+\cdots+Z_3x_3)\bigr)^{\otimes s}
       \otimes S_\eta(B)v_\eta.
\]
Here the pencil coefficients, coefficients of the linear form $L$,
and entries of $B\in\operatorname{End}(V^*)$ are independent.
The first factor has parameter degree $ns$, exactly as for the
determinant parametrization. Therefore the projective image of
the $E$th invariant Veronese has degree at most
$(pE)^{v-1}\leq Q_P$. The same Chardin upper bound, enlarged to
dimension $\bar v$, and Noether-normalization lower bound give
\[
 \dim(A_P)_{E(k+2)}\leq Q_P\binom{k+\bar v+1}{\bar v-1},
 \qquad
 \dim(A_D)_{E(k-1)}\geq\binom{k+\bar v-1}{\bar v}.
\]
Use \eqref{G18:sandwich} with the roles reversed and the preceding
binomial calculation with $\bar v$ in place of $u$.
This proves \eqref{G18:everylargereverse}. If
$n\geq\sqrt2\,m$, then
$2n^2-(4m^2-2m+3)\geq2m-3>0$, as required.

\paragraph{4. The all-integer asymptotic ratio.}
For a standard graded domain $B$ of dimension $q\geq2$, write
$e(B)=\deg\operatorname{Proj}B>0$. Its Hilbert function satisfies
\[
 \dim B_k=\frac{e(B)}{(q-1)!}k^{q-1}+O(k^{q-2}).
\]
This is the usual Hilbert-polynomial normalization used in G-10.
One can check the leading term directly after linear Noether
normalization: $B$ is a finite torsion-free module over a polynomial
ring in $q$ degree-one variables; a homogeneous basis over its
fraction field spans a free submodule of rank $e(B)$, and the
torsion quotient has growth $O(k^{q-2})$. An annihilating homogeneous
polynomial for that quotient and its finite generating set give
this error bound from the polynomial-ring Hilbert function.

Apply this formula to $B_D=A_D^{(E)}$ and $B_P=A_P^{(E)}$.
Their dimensions are $u,v$, both at least two here. Indeed
$c_m\geq20$: R-13 gives $c_3=20$, and a block diagonal pencil
with a $3\times3$ block and $m-3$ entries $x_0$ embeds
$x_0^{m-3}\mathcal P_{3,4}$ in $\mathcal P_{m,4}$.
Also $r_\eta\geq0$, including the rectangular case $\eta=0$.
Sandwich
\eqref{G18:sandwich}, uniformly over its finitely many residues,
gives for $A$ of dimension $q$
\[
 \dim A_j=\frac{e(A^{(E)})}{E^{q-1}(q-1)!}
                 j^{q-1}+O(j^{q-2}).
\]
Taking the quotient of these two positive leading terms proves
\eqref{G18:ratio}, with
\[
 \kappa_\nu=
 \frac{e(B_P)}{e(B_D)}\,
 \frac{(u-1)!}{(v-1)!}\,E^{u-v}>0,
 \qquad v-u=c_m-2n^2+1.
\]
All multiplicities are identified with the full orbit-closure
ones by the same four-variable highest-weight restriction as in
G-17. The final occurrence assertion follows from
\eqref{G18:occurrence} in the strict region. Outside that region,
G-05 and G-14 give the opposite weak multiplicity inequality at
every stretch, so an occurrence obstruction is already impossible.
\end{proof}

\paragraph{The precise remaining range.}
The reverse inequality concerns sufficiently large stretches of
each fixed direction. It does not exclude an individual
finite-degree gap, give a threshold uniform over changing
directions, or prove containment of orbit closures. Rectangular
directions obey the padding boundary of G-12 as well as the
all-stretch conclusion in the strict range. When
$c_m+3=2n^2+2$, the leading quotient degrees, and
possibly lower Hilbert coefficients, are the remaining comparison.

